\section*{What the released adapter determines}

This note states and proves what a released LoRA pair $(A_T,B_T)$ determines about the data it was
trained on. It is theorem-first: each statement is proved, or marked as measured with the job that
measured it, or marked open. Nothing here is inferred from a count alone.

\paragraph{Setting.} One adapted layer $\ell$ with input dimension $d_\ell$, output dimension
$m_\ell$, LoRA rank $r_\ell$; the frozen weight is $W_\ell^0$ and the adapter makes it
$W_\ell^0+B_\ell A_\ell$. Let $H_\ell=[h_{\ell,1},\dots,h_{\ell,N}]\in\mathbb R^{d_\ell\times N}$ be
the representations that entered the layer during training, and
\[
  q_\ell \;:=\; \operatorname{rank}H_\ell .
\]
Throughout, \textbf{$q$ and not $N$ is the operative quantity}: the release determines the dimension
of the span of the training representations, never the number of photographs. $q_\ell=N$ only when
the $N$ representations are linearly independent, and augmentation makes $N_{\text{repr}}\ge q_\ell$
with $N_{\text{image}}$ smaller still. Training is $B_0=0$, $A_0$ drawn from a density, plain
gradient descent, and the layer's inputs are fixed across steps.

\medskip
Write $S_\ell=\operatorname{row}(B_{\ell,T})\subseteq\mathbb R^{r_\ell}$, let $\Pi_\ell$ be the
orthogonal projector onto $S_\ell^{\perp}$, and define the \emph{certificate}
\[
  C_\ell \;:=\; \Pi_\ell\,A_{\ell,T}, \qquad s_\ell:=\operatorname{rank}C_\ell .
\]
$C_\ell$ is computed from the released factors alone --- no recipe, no labels, no seed, no $N$.

\subsection*{1. The certificate annihilates the training span}

\begin{theorem}\label{thm:cert}
Assume
\begin{itemize}
\item[(H1)] $B_0=0$, plain gradient descent with no weight decay, and the layer's inputs fixed
  across steps;
\item[(H2)] $\operatorname{rank}A_0=r$ and $\operatorname{col}(H)\cap\ker A_0=\{0\}$ --- generic
  over the seed, and needed because $A_0$ maps $\mathbb R^{d}$ to $\mathbb R^{r}$ with $r<d$ and so
  has a kernel of dimension $d-r$;
\item[(H3)] \emph{excitation}: $\operatorname{rank}B_T=q$. Since
  $\operatorname{row}(B_T)\subseteq\operatorname{col}(A_0H)$ always, $\operatorname{rank}B_T\le q$
  unconditionally, and (H3) is the case of equality. The attacker can compute
  $\operatorname{rank}B_T$ from the release and so holds a \emph{lower bound} on $q$; verifying that
  the bound is tight needs $H$ and is not available to them. So (H3) is a hypothesis about the
  release, not a test the attacker can run --- but it is the one whose failure has a visible
  direction, since a collapsed $\operatorname{rank}B_T$ means the certificate is built from too small
  a subspace and $CH\neq0$.
\end{itemize}
Nothing here requires $I+M_tG$ to be invertible. That condition is one route by which (H3) comes to
hold, but it is not used in the proof, and (H3) is checked directly rather than inferred.
Then
\[
  \Pi A_T=\Pi A_0\ \text{ exactly},\qquad
  C\,H=0,\qquad s=r-q,\qquad \ker C=\operatorname{span}(H)\oplus\ker A_0 .
\]
\end{theorem}

\begin{proof}
Under (H1) the factor dynamics close: $B_t=P_t(A_0H)^{\top}$ and $A_t=A_0(I+HM_tH^{\top})$ for
coefficient matrices $P_t,M_t$. The induction needs only \emph{containment}: $A_tH=A_0H\,(I+M_tG)$
lies in $\operatorname{col}(A_0H)$ whatever the right factor is, so the $B$ update keeps its rows
there and the $A$ update keeps its shape, with no condition on $I+M_tG$.

Every row of $B_T=P_T(A_0H)^{\top}$ is a combination of the columns of $A_0H$, so
$S=\operatorname{row}(B_T)\subseteq\operatorname{col}(A_0H)$ always. Under (H2),
$\dim\operatorname{col}(A_0H)=\operatorname{rank}(A_0H)=q$, and (H3) gives $\dim S=q$; a subspace
contained in another of the same finite dimension equals it, so
$S=\operatorname{col}(A_0H)$. \emph{Excitation is therefore a consequence of (H3), not an extra
assumption about the training signal.}

Now $A_T=A_0+A_0HM_TH^{\top}$, and $\Pi$ annihilates $\operatorname{col}(A_0H)=S$, so
$\Pi A_0H=0$ and the second term dies:
\[
  \Pi A_T=\Pi A_0+(\Pi A_0H)M_TH^{\top}=\Pi A_0 ,
\]
an exact identity with no scalar. Consequently $C=\Pi A_T=\Pi A_0$, and
$CH=\Pi A_0H=0$.

For the rank, $\Pi$ has rank $r-\dim S=r-q$ by (H3); $A_0$ is onto $\mathbb R^{r}$ by (H2), so
$\Pi A_0$ has the same rank as $\Pi$ and $s=r-q$. Finally $Cx=0$
iff $A_0x\in\operatorname{col}(A_0H)$ iff $x\in\operatorname{span}(H)+\ker A_0$, and (H2) makes the
sum direct.
\end{proof}

\begin{remark}[the certificate is a rejection test]\label{rem:rejection}
$Ch=0$ is exact for \emph{subspace} consistency and nothing more. Every point of
$\operatorname{span}(H)$ passes it, including points that are not training representations. It can
refute a candidate; it cannot accept one. This is used again in Lemma~\ref{lem:hull}.
\end{remark}

\begin{remark}[homogeneity]
$C(\alpha h)=\alpha Ch$, so a candidate's norm is unconstrained: the certificate cannot determine
$\|h\|$. Whether replay determines it is an experiment, not a theorem.
\end{remark}

\subsection*{2. Quotient sensing: what the release reveals about the seed}

\begin{proposition}\label{prop:sensing}
Under the hypotheses of Theorem~\ref{thm:cert}, $C=\Pi A_0$ --- \emph{with no free scalar}.
Consequently, for a fixed $h$ with off-span component $h_\perp$ and Gaussian $A_0$ with entry scale
$\sigma_0$,
\[
  \frac{\|Ch\|^{2}}{\sigma_0^{2}\|h_\perp\|^{2}}\;\sim\;\chi^{2}_{\,r-q},
\]
while every $h\in\operatorname{span}(H)$ gives exactly zero.
\end{proposition}

\begin{proof}
Immediate from $\Pi A_T=\Pi A_0$ in Theorem~\ref{thm:cert}. $\Pi A_0$ restricted to
$\operatorname{span}(H)^{\perp}$ is a Gaussian matrix of rank $r-q$ acting on $h_\perp$, giving the
stated $\chi^{2}$.
\end{proof}

Three consequences follow, and the first is sharper than the form this project carried until now.
The release pins the $S^{\perp}$-component of the seed \emph{exactly}, not up to a scalar, so the
seed reduction leaves $q\,d$ unknowns rather than $q\,d+1$. The constant is not fitted: it is $1$ by
the proof, and measured at $1.000000000000$ on the gate of job 675031. And a measured $c_T\neq1$ is
then a \emph{diagnostic} rather than noise --- it is the signature of weight decay or of an optimiser
that is not plain gradient descent, either of which breaks (H1).

This gives a membership test that needs no shadow models, no recipe and no distributional
assumption --- an assumptions result rather than a detection result.

\subsection*{3. What a chart of dimension $k$ leaves free}

An attacker searches a chart $\psi:\mathbb R^{k}\to\mathbb R^{d}$ rather than all of $\mathbb R^d$.

\begin{theorem}[generic tangent rank, one adapted layer]\label{thm:tangent}
For a single adapted layer, with $U_x$ an orthonormal basis of the chart's tangent space at $x$ and
$d_\parallel$ the number of tangent directions lying inside $\operatorname{span}(H)$,
\[
  \operatorname{rank}\bigl(C\,J_\psi(x)\,U_x\bigr)\;=\;\min\bigl(r-q,\;k-d_\parallel\bigr)
  \quad\text{almost surely over the seed.}
\]
Cross-layer stacking is \emph{not} covered by this statement; the rank of the stacked Jacobian across
adapted layers is unmeasured.
\end{theorem}

\subsection*{4. The certificate can never separate the span --- and what fixes it}

\begin{lemma}[affine-hull degeneracy]\label{lem:hull}
$C$ is linear and $CH=0$, so $C\bigl(\sum_i\alpha_ih_i\bigr)=0$ for \emph{every} coefficient vector
$\alpha$. Hence the certificate alone cannot distinguish the training representations from their
span. Moreover, let $\Psi=\varphi\circ\psi$ be the composition from chart coordinates to the adapted
layer's input, with the privates on-chart, $h_i=\Psi(z_i)$. If $\Psi$ is affine then for every
$\alpha$ with $\sum_i\alpha_i=1$,
\[
  \Psi\Bigl(\sum_i\alpha_iz_i\Bigr)=\sum_i\alpha_ih_i\in\operatorname{col}H\subseteq\ker C ,
\]
so the certificate's zero set on the chart contains an affine subspace of dimension $\min(N-1,k)$
through the training points, and they are not isolated zeros at any $k$.
\end{lemma}

\begin{proof}
The first claim is linearity. For the second, $\sum_i\alpha_i=1$ is exactly what absorbs the
intercept: for affine $\Psi(z)=Lz+b$, $\Psi(\sum\alpha_iz_i)=L\sum\alpha_iz_i+b=\sum\alpha_i(Lz_i+b)$.
The affine hull of $N$ points of $\mathbb R^k$ has dimension $\min(N-1,k)$ when the $z_i$ are
affinely independent --- generic --- and contains each $z_i$.
\end{proof}

\begin{remark}[why a count cannot see this]\label{rem:count-blind}
The capacity count is a transversality count and is valid off $\operatorname{col}H$; it is invalid
\emph{on} it, because $\operatorname{col}H\subseteq\ker C$ is forced by construction rather than
generic. In the language of Theorem~\ref{thm:tangent} this is $d_\parallel>0$. Nonlinearity between
chart and adapter is how the condition is bought, not what it says; the condition is that the chart
meet the training span only at the training points.
\end{remark}

\begin{remark}[the available checks cannot substitute for it]\label{rem:checks}
$\|CH\|$, $\operatorname{rank}C$, the excitation gap and the residual at the truths are functions of
$C$ and $H$ alone, and none of them mentions the chart. Holding $C$ and $H$ fixed while varying the
chart moves the training points from isolated to non-isolated with every one of those numbers
unchanged. \emph{No} function of $(C,H)$ alone can decide the condition of
Remark~\ref{rem:count-blind}. One test does bear on it and is one-sided: at a found $z^{\ast}$,
$J=C\,D\Psi(z^{\ast})$ of full column rank $k$ certifies the zero is isolated, while
$\operatorname{rank}J<k$ certifies nothing, a degenerate isolated zero having it too.
\end{remark}

\paragraph{Measured (job 331384; register: read-rows).} One release, one affine chart
($\psi$ linear, $\varphi$ the identity), the same random starts, with $k=12$, $N=8$, $r=24$, $m=20$
chosen below both capacity lines and above $N-1$ so that the hypothesis of Lemma~\ref{lem:hull} is
the only one failing. The certificate route recovered nothing from 60 starts; every returned point
was an exact affine combination of the training representations with coefficient sum
$1.000000000$ at \emph{every one} of the 60 starts (maximum deviation $8.9\times10^{-16}$), and the
isolation rank was $5$,
a deficit of exactly $N-1=7$. Replay on the same release and the same starts returned all eight
representations from 18 of those 60 starts with a worst-image error at or below
$2.2\times10^{-15}$ and a replay residual below $10^{-12}$; a 19th start cleared the $10^{-2}$ landing
bar at a worst-image error of $4.4\times10^{-3}$. Zero aliases at either bar.

\subsection*{5. Replay resolves what the certificate cannot}

The two routes have nested zero sets, in one direction only:
\[
  \{\text{truth}\}\;\subseteq\;\{\rho=0\}\;\subseteq\;\{CH=0\},
\]
where $\rho$ is the residual of simulating the training recipe on candidate data against the released
pair. So $\rho=0\Rightarrow CH=0$ and never the converse: a blend is an exact zero of $C$ and is not
generally a zero of $\rho$, because reproducing $B_T$ means reproducing what the recipe actually did.
The measurement above is this inclusion being strict.

\begin{remark}[identifiability is not a property of the release alone --- it depends on chart and route]
The same chart that makes one channel provably blind leaves the other exact. Identifiability is
therefore not a property of the release alone, nor of the (release, chart) pair, but of all three.
A corollary: an adapter on the input layer read through a linear chart is non-identifying
\emph{for the certificate route}; the corresponding statement for replay is false.
\end{remark}

\subsection*{6. The count, as arithmetic}

Three quantities are easy to conflate, and only the third bears on identifiability.

\begin{enumerate}
\item \textbf{Directions.} $\sum_\ell s_\ell$, with $s_\ell=r_\ell-q_\ell$ --- constraint directions
      per image.
\item \textbf{Raw scalar equations on the training set.} $C_\ell H_\ell=0$ is an $s_\ell\times N$
      system, so each adapted layer contributes $s_\ell N$ \emph{raw scalar equalities} and the release
      contributes $N\sum_\ell s_\ell$ raw rows in total. This is a row count, not information.
\item \textbf{Independent information on a chart.} The rank of the stacked Jacobian $J_F$, where
      $F(z)$ stacks $C_\ell\Phi^0_\ell(G(z))$ over the adapted layers. For one layer this is
      Theorem~\ref{thm:tangent}; across layers it is \emph{unmeasured}.
\end{enumerate}

Only the third bears on identifiability, and the chain from the first to a defensible statement has
four links, each of which can lose information: \textbf{raw rows} $\to$ \textbf{per-example
restrictions} $\to$ \textbf{rank on the chart actually searched} $\to$ \textbf{remaining joint
degrees of freedom}. The third link depends on which chart regime the attacker is in --- a
\emph{universal} chart over all natural images, a \emph{target-conditioned} chart built from public
data about a declared target, or a \emph{shared-concept} chart $x_i=G(z_{\text{shared}},z_i)$ ---
and the three differ in the $k$ needed to represent the private data at a given fidelity. That
comparison is being measured; it is not assumed here.

Worked from the architecture: with $N=8$, $r=16$, $q=8$, each module gives $s=8$, so one module is
$8$ directions per image and $64$ raw equalities; four modules $32$ and $256$; eight modules $64$ and
$512$; twelve modules $96$ and $768$. With $r=32$ and $q=8$, $s=24$, and eight modules give $192$ per
image and $1536$ raw. Deep modules inherit the caveat that exactness holds only where the layer's
inputs are fixed during training.

\begin{remark}[multiplicity]
If images are parameterised independently, the system is block-diagonal: each image's
$\sum_\ell s_\ell$ equations act on its own $k$ unknowns, ranks add across images but never exceed
$k$ per image, and multiplying both sides by $N$ changes no ratio. The relevant comparison is
$\sum_\ell s_\ell$ per image against $k$ per image. Only a chart that \emph{shares} parameters,
$x_i=G(z_{\text{shared}},z_i)$, pools all $N\sum_\ell s_\ell$ equations onto the shared block.
Replay differs: all images are coupled through the common released endpoint even when their latents
are independent, so its identifiability is joint from the start.
\end{remark}

\begin{remark}[what a count may and may not conclude]\label{rem:count-language}
A configuration is \emph{ruled out by count}, \emph{not ruled out by count}, or \emph{marginal ---
requires a measured Jacobian rank}. ``Not ruled out by count'' never means identifiable:
identifiability additionally requires that the chart meet the training span only at the training
points, and the count does not test that condition (Lemma~\ref{lem:hull},
Remark~\ref{rem:count-blind}). No reconstruction claim follows from a count alone.
\end{remark}

\subsection*{7. Status}

Proved here: Theorem~\ref{thm:cert}, Proposition~\ref{prop:sensing}, Lemma~\ref{lem:hull} and the
remarks. Theorem~\ref{thm:tangent} is proved for a single layer; the stacked-Jacobian rank across
layers is open and is the quantity item~3 of \S6 actually needs. The capacity line $k<m+r-q$ for the
reduced $B_T$ channel is measured sharp to one unit of $k$ on synthetic and MNIST releases and
remains a candidate for the full factor pair. The deep-layer perturbative bound is not proved and is
stated as open; a counterexample would be written up as one.
